% Generated draft: requires mathematical and source review.
\documentclass[11pt]{article}
\usepackage[margin=25mm]{geometry}
\usepackage{fontspec}
\setmainfont{Latin Modern Roman}
\usepackage{amsmath,amssymb,mathtools,unicode-math}
\setmathfont{Latin Modern Math}
\usepackage{enumitem,needspace,xurl,hyperref}
\setlength{\parindent}{0pt}
\setlength{\parskip}{5pt}
\setlength{\emergencystretch}{3em}
\newcommand{\N}{\mathbb{N}}
\newcommand{\Z}{\mathbb{Z}}
\newcommand{\Q}{\mathbb{Q}}
\newcommand{\R}{\mathbb{R}}
\newcommand{\C}{\mathbb{C}}
\newcommand{\F}{\mathbb{F}}
\newcommand{\A}{\mathbb{A}}
\newcommand{\PP}{\mathbb{P}}
\newcommand{\E}{\mathbb{E}}
\newcommand{\eps}{\varepsilon}
\newcommand{\dd}{\,\mathrm{d}}
\begin{document}
\section*{UnsolvedMath ID: 30004694}
\textbf{Source problem code:} catalogue entry
\subsection{Definitions and mathematical statement}
The supplied record poses a historical and interpretive question rather than a conjecture with specified mathematical objects, hypotheses, and a truth-valued conclusion. It supplies no formal definitions of mathematical ``classics,'' ``progress,'' or ``Antiquity,'' and specifies no corpus, geographical scope, historical boundaries, or criterion for an adequate explanation. These terms therefore remain unresolved source terminology.

The source question is:
\begin{quote}
Why does mathematics, a discipline that values progress, have classics, especially classics from Antiquity?
\end{quote}
The characterization of mathematics as valuing progress is part of the question's framing. The request concerns an explanation of the persistence of classics, with particular emphasis on ancient classics. A precise mathematical proposition cannot be extracted from this record without introducing unsupported definitions or changing the intended historical inquiry.
\subsection{Short English statement}
Explain why mathematical classics persist in a discipline described as valuing progress, especially why ancient works retain this role. The source leaves the relevant terms, historical scope, and standards of explanation unspecified; this draft preserves the historical question without inventing a mathematical conjecture.
\subsection{Sources}
\begin{itemize}
\item Local UnsolvedMath registry, permanent ID \(30004694\), title ``The Persistence of Ancient Mathematical Classics,'' exact source code \texttt{catalogue entry}; supplied complete source record from the Oberwolfach Reports Open Problems collection. Repository registry locations are \texttt{data/unsolvedmath/problems.json} and the supplementary \texttt{catalog/problems.json}; the supplied record does not identify which contains this entry.
\item Supplied source citation: \emph{Mathematics and its Ancient Classics Worldwide: Translations, Appropriations, Reconstructions, Roles} (2022), Oberwolfach Reports, pp.\ 1347--1406, DOI \texttt{10.4171/owr/2021/26}. Supplied source and literature link: \url{https://doi.org/10.4171/owr/2021/26}. The background describes Oberwolfach Reports as documenting mathematical workshops and their open-problem discussions. No individual proposer is supplied; the registry gives proposed year 2021.
\item Supplied literature assessment, checked on \texttt{2026-08-22}: the record labels the question ``open'' and characterizes it as a broad historical inquiry for which no unique, settled answer is established in the primary literature. The listed checked source is the Oberwolfach report linked above. This is supplied dated curation, not an independently verified assessment; the record explicitly requests independent verification of time-sensitive status and cited claims.
\item The supplied statement status is \texttt{corrected\_verified}. Its accompanying verification note reports a manual semantic, typographic, and LaTeX review confirming that the clean statement is self-contained and preserves the target without unsupported additions. This is a supplied review claim, not independent verification or evidence that the inquiry admits a precise mathematical formulation.
\end{itemize}
\end{document}
